\chapter{来源、固定版本与插图出处}

\label{seedance-sources}

本章从来源仓库导入登记信息，用于追溯正文的基础方法和候选设计。论文结论受原实验设置限制，引用不代表本文组合方案已经验证。网页状态沿用2026年11月22日记录；本次未刷新线上产品与价格。各模块仍为审阅稿。

\section{论文与网页}

\sourceentry{adaspa}{Training-free and Adaptive Sparse Attention for Efficient Long Video Generation}{arXiv 2502.21079v1; arXiv 2502.21079v1; 来源登记检查日 2026-11-19}{https://arxiv.org/abs/2502.21079v1}

\sourceentry{alchemint}{AlcheMinT: Fine-grained Temporal Control for Multi-Reference Consistent Video Generation}{CVPR 2026 publisher edition; arXiv 2512.10943; 来源登记检查日 2026-11-22}{https://openaccess.thecvf.com/content/CVPR2026/html/Girish_AlcheMinT_Fine-grained_Temporal_Control_for_Multi-Reference_Consistent_Video_Generation_CVPR_2026_paper.html}

\sourceentry{dancegrpo}{DanceGRPO: Unleashing GRPO on Visual Generation}{arXiv 2505.07818v4; arXiv 2505.07818v4; 来源登记检查日 2026-11-22}{https://arxiv.org/abs/2505.07818v4}

\sourceentry{dit}{Scalable Diffusion Models with Transformers}{arXiv 2212.09748v2; arXiv 2212.09748v2; 来源登记检查日 2026-11-19}{https://arxiv.org/abs/2212.09748v2}

\sourceentry{geometry-as-context}{Geometry-as-context: Modulating Explicit 3D in Scene-consistent Video Generation to Geometry Context}{CVPR 2026 publisher edition; arXiv 2602.21929; 来源登记检查日 2026-11-22}{https://openaccess.thecvf.com/content/CVPR2026/html/Hu_Geometry-as-context_Modulating_Explicit_3D_in_Scene-consistent_Video_Generation_to_Geometry_CVPR_2026_paper.html}

\sourceentry{long-context-tuning}{Long Context Tuning for Video Generation}{arXiv 2503.10589v1; arXiv 2503.10589v1; 来源登记检查日 2026-11-22}{https://arxiv.org/abs/2503.10589v1}

\sourceentry{longlive}{LongLive: Real-time Interactive Long Video Generation}{arXiv 2509.22622v2; arXiv 2509.22622v2; 来源登记检查日 2026-11-22}{https://arxiv.org/abs/2509.22622v2}

\sourceentry{ltx-2}{LTX-2: Efficient Joint Audio-Visual Foundation Model}{arXiv 2601.03233v1; arXiv 2601.03233v1; 来源登记检查日 2026-11-22}{https://arxiv.org/abs/2601.03233v1}

\sourceentry{ltx-original-model}{Lightricks/LTX-2 original dev model metadata}{official\_model\_repository; 来源登记检查日 2026-11-22}{https://huggingface.co/Lightricks/LTX-2/tree/64f65d09e0b76c96202b5bc09997673184cf9b5d}

{\small 支持范围：Original dev identity, repository configs and 603120-byte native safetensors header read; tensor payload not downloaded; remote LFS hash is not a local full-file check. Detailed metadata: research/2026-11-22-implementation-contracts/model-config-check.json\par}

\sourceentry{magvit-v2}{Language Model Beats Diffusion — Tokenizer is Key to Visual Generation}{arXiv 2310.05737v3; arXiv 2310.05737v3; 来源登记检查日 2026-11-22}{https://arxiv.org/abs/2310.05737v3}

\sourceentry{phantom}{Phantom: Subject-consistent video generation via cross-modal alignment}{arXiv 2502.11079v2; arXiv 2502.11079v2; 来源登记检查日 2026-11-22}{https://arxiv.org/abs/2502.11079v2}

\sourceentry{rayflow}{RayFlow: Instance-Aware Diffusion Acceleration via Adaptive Flow Trajectories}{arXiv 2503.07699v2; arXiv 2503.07699v2; 来源登记检查日 2026-11-22}{https://arxiv.org/abs/2503.07699v2}

\sourceentry{rewarddance}{RewardDance: Reward Scaling in Visual Generation}{arXiv 2509.08826v1; arXiv 2509.08826v1; 来源登记检查日 2026-11-22}{https://arxiv.org/abs/2509.08826v1}

\sourceentry{rope}{RoFormer: Enhanced Transformer with Rotary Position Embedding}{arXiv 2104.09864v1; arXiv 2104.09864v1; 来源登记检查日 2026-11-22}{https://arxiv.org/abs/2104.09864v1}

\sourceentry{seaweed-apt}{Diffusion Adversarial Post-Training for One-Step Video Generation}{arXiv 2501.08316v3; arXiv 2501.08316v3; 来源登记检查日 2026-11-19}{https://arxiv.org/abs/2501.08316v3}

\sourceentry{seedance-1.0}{Seedance 1.0: Exploring the Boundaries of Video Generation Models}{arXiv 2506.09113v2; arXiv 2506.09113v2; 来源登记检查日 2026-11-22}{https://arxiv.org/abs/2506.09113v2}

\sourceentry{seedance-1.5-pro}{Seedance 1.5 pro: A Native Audio-Visual Joint Generation Foundation Model}{arXiv 2512.13507v3; arXiv 2512.13507v3; 来源登记检查日 2026-11-22}{https://arxiv.org/abs/2512.13507v3}

\sourceentry{seedance-2.0}{Seedance 2.0: Advancing Video Generation for World Complexity}{arXiv, 2026-04-15; model card / technical report; arXiv 2604.14148v1; 来源登记检查日 2026-11-19}{https://arxiv.org/abs/2604.14148v1}

\sourceentry{seedance-2.5-product}{Seedance 2.5 product page}{official\_product; 来源登记检查日 2026-11-22}{https://seed.bytedance.com/en/seedance2_5}

{\small 支持范围：动态英文页；30秒、延长两次、参考/编辑/白模/绿幕。与发布页“多轮延长”分别保留范围。\par}

\sourceentry{seedance-2.5-release}{一镜成片，随心参考｜Seedance 2.5 正式发布}{official\_release; 来源登记检查日 2026-11-22}{https://seed.bytedance.com/blog/one-take-creation-flexible-referencing-introducing-seedance-2-5}

{\small 支持范围：实际访问并阅读正文；支持产品能力和架构连续性陈述，不披露完整实现。API近期上线是发布时表述，未核查当前API文档。\par}

\sourceentry{tarsier2}{Tarsier2: Advancing Large Vision-Language Models from Detailed Video Description to Comprehensive Video Understanding}{arXiv 2501.07888v1; arXiv 2501.07888v1; 来源登记检查日 2026-11-22}{https://arxiv.org/abs/2501.07888v1}

\sourceentry{vace}{VACE: All-in-One Video Creation and Editing}{arXiv 2503.07598v2; arXiv 2503.07598v2; 来源登记检查日 2026-11-22}{https://arxiv.org/abs/2503.07598v2}

\sourceentry{video-dit-scaling-laws}{Towards Precise Scaling Laws for Video Diffusion Transformers}{arXiv 2411.17470v2; arXiv 2411.17470v2; 来源登记检查日 2026-11-19}{https://arxiv.org/abs/2411.17470v2}

\sourceentry{videocof}{VideoCoF: Unified Video Editing with Temporal Reasoner}{CVPR 2026 publisher edition; arXiv 2512.07469; 来源登记检查日 2026-11-22}{https://openaccess.thecvf.com/content/CVPR2026/html/Yang_VideoCoF_Unified_Video_Editing_with_Temporal_Reasoner_CVPR_2026_paper.html}

\sourceentry{vincie}{VINCIE: Unlocking In-context Image Editing from Video}{arXiv 2506.10941v2; arXiv 2506.10941v2; 来源登记检查日 2026-11-19}{https://arxiv.org/abs/2506.10941v2}

\sourceentry{vincie-next}{VINCIE-NExT: Unlocking Video Editing from Images via In-Context Modeling}{NeurIPS 2026 (according to the arXiv authors' comment); arXiv 2610.12104v1; 来源登记检查日 2026-11-19}{https://arxiv.org/abs/2610.12104v1}

\sourceentry{wan-2.1}{Wan: Open and Advanced Large-Scale Video Generative Models}{arXiv 2503.20314v2; arXiv 2503.20314v2; 来源登记检查日 2026-11-22}{https://arxiv.org/abs/2503.20314v2}

\section{上游代码与模型身份}

\sourceentry{dancegrpo-code}{XueZeyue/DanceGRPO}{固定 commit 15cc71d53cc2e6e18a68ee607d5fb6ba9a99e344；代码许可 Apache-2.0}{https://github.com/XueZeyue/DanceGRPO/tree/15cc71d53cc2e6e18a68ee607d5fb6ba9a99e344}

{\small 支持范围：固定版 README 将 Wan-2.1 路线当作图像生成并用 HPS 奖励，不能当成已验证的完整视频对齐配方。\par}

\sourceentry{diffsynth-code}{modelscope/DiffSynth-Studio}{固定 commit d32f72424cd22bc0ea8100cc57559149b4eac933；代码许可 Apache-2.0}{https://github.com/modelscope/DiffSynth-Studio/tree/d32f72424cd22bc0ea8100cc57559149b4eac933}

{\small 支持范围：真实训练框架；现成 LoRA 不自动提供多参考绑定或局部编辑能力。\par}

\sourceentry{longlive-code}{NVlabs/LongLive}{固定 commit 31c70d9071cd7851d82939806dc9fac645fcabe5；代码许可 Apache-2.0}{https://github.com/NVlabs/LongLive/tree/31c70d9071cd7851d82939806dc9fac645fcabe5}

{\small 支持范围：固定阅读 LongLive1.0 子目录；仓库根目录另含2.0/Plug/Long-WAM，不能混用入口。\par}

\sourceentry{ltx-code}{Lightricks/LTX-2}{固定 commit daa6b16fc7f0f1628f2c9e36a303e3312b333cc6；代码许可 Apache-2.0}{https://github.com/Lightricks/LTX-2/tree/daa6b16fc7f0f1628f2c9e36a303e3312b333cc6}

{\small 支持范围：当前固定版本同时覆盖多个 LTX 型号，模型/codec/LoRA 必须匹配；代码 Apache-2.0，模型另有 Community License。\par}

\sourceentry{ltx-original-code}{Lightricks/LTX-2}{固定 commit bd92a5f408eb8c99f8c05e9fe2fcaaec8b56b638；代码许可 LTX-2 Community License (root LICENSE at this historical commit)}{https://github.com/Lightricks/LTX-2/tree/bd92a5f408eb8c99f8c05e9fe2fcaaec8b56b638}

{\small 支持范围：固定2026-01原系列版本；不能与后来LTX2.5的checkpoint、文本和codec混用。该历史commit根LICENSE为LTX-2 Community License，不能套用当前版本Apache代码许可。\par}

\sourceentry{phantom-code}{Phantom-video/Phantom}{固定 commit bd84b602dcc949e23c89cbbf266b6f5975f2f025；代码许可 Apache-2.0}{https://github.com/Phantom-video/Phantom/tree/bd84b602dcc949e23c89cbbf266b6f5975f2f025}

{\small 支持范围：发布的是 Phantom-Wan 适配，不能逐项当成论文中内部基座；该版本 README 的训练代码仍为未勾选。\par}

\sourceentry{tarsier-code}{bytedance/tarsier}{固定 commit 63de87a760c2aec6b31470a08046969a13787126；代码许可 Apache-2.0}{https://github.com/bytedance/tarsier/tree/63de87a760c2aec6b31470a08046969a13787126}

{\small 支持范围：区分 Tarsier/Tarsier2 checkpoint 和其 processor/config；自动描述仍需事实抽检。\par}

\sourceentry{vace-code}{ali-vilab/VACE}{固定 commit 48eb44f1c4be87cc65a98bff985a26976841e9f3；代码许可 Apache-2.0}{https://github.com/ali-vilab/VACE/tree/48eb44f1c4be87cc65a98bff985a26976841e9f3}

{\small 支持范围：VACE不同基座分支不自动兼容本方案LTX-2；不是其现成插件。\par}

\sourceentry{videocof-code}{knightyxp/VideoCoF}{固定 commit f45b8009616b3ef8e2024a7da585249b4c74fc2b；代码许可 Apache-2.0}{https://github.com/knightyxp/VideoCoF/tree/f45b8009616b3ef8e2024a7da585249b4c74fc2b}

{\small 支持范围：Wan/VideoX-Fun工程；论文联合区域/目标去噪不等于本文显式两阶段实现。\par}

\sourceentry{vincie-code}{ByteDance-Seed/VINCIE}{固定 commit 97cd2c5701e694dabcfea320c0cf59893d427fab；代码许可 Apache-2.0}{https://github.com/ByteDance-Seed/VINCIE/tree/97cd2c5701e694dabcfea320c0cf59893d427fab}

{\small 支持范围：主要用于多轮图像编辑；不能把此入口称为时间戳视频编辑或直接替代 VINCIE-NExT。\par}

\sourceentry{wan21-code}{Wan-Video/Wan2.1}{固定 commit 9737cba9c1c3c4d04b33fcad41c111989865d315；代码许可 Apache-2.0}{https://github.com/Wan-Video/Wan2.1/tree/9737cba9c1c3c4d04b33fcad41c111989865d315}

{\small 支持范围：基础推理入口；不提供本仓库候选 reference adapter 的训练方案。\par}

\section{论文插图}

插图是原论文的局部示意，不能当作本书候选系统或 Seedance 2.5 官方架构。以下页序从1开始；裁剪坐标以 PDF point 为单位。原 PDF 哈希与完整坐标随附于插图 manifest，未复制整篇论文。

\begin{longtable}{@{}p{0.35\linewidth}p{0.25\linewidth}p{0.30\linewidth}@{}}

\toprule 插图 & 来源 ID & PDF页序与裁剪范围 \\ \midrule\endhead

audio-video-pipeline & seedance-1.5-pro & 2; [65, 522, 545, 673] \\

causal-tokenizer & magvit-v2 & 5; [105, 57, 507, 225] \\

data-pipeline & seedance-1.0 & 6; [65, 70, 546, 214] \\

edit-data & vincie & 4; [105, 65, 508, 264] \\

generative-reward & rewarddance & 5; [65, 65, 546, 368] \\

geometry-context & geometry-as-context & 5; [317, 357, 555, 626] \\

joint-av-block & ltx-2 & 5; [106, 55, 506, 243] \\

multishot-positions & long-context-tuning & 5; [57, 55, 555, 199] \\

reason-edit & videocof & 3; [57, 55, 556, 276] \\

spatial-temporal-dit & seedance-1.0 & 5; [65, 70, 547, 278] \\

streaming-cache & longlive & 2; [105, 85, 507, 269] \\

subject-reference & phantom & 4; [57, 60, 555, 246] \\

timed-reference & alchemint & 4; [57, 57, 556, 252] \\

training-pipeline & seedance-1.0 & 8; [65, 480, 546, 707] \\

unified-context & vace & 4; [57, 55, 556, 211] \\

wan-block & wan-2.1 & 14; [345, 130, 506, 333] \\

\bottomrule\end{longtable}

\section{本次导入身份}

来源仓库：\url{https://github.com/foocker/multimodal-journey}。

固定 commit：\codepath{0f1b6bb86d609d75d8c2582063d85707d0180d39}。导入日：2026年11月24日。

逐模块源文件 SHA-256、图件来源、原模型元数据与新旧文件对应关系，见随附的 \codepath{research/seedance-import/manifest.json}。这份记录是文件追溯，不能代替科学复现。
